\chapter{Conclusões e Trabalhos Futuros} \label{Cap:Conclusoes}

\label{chp:conclusao}

\section{Considerações iniciais}
\label{sec:concl_consideracoes_iniciais}

Este capítulo encerra a monografia apresentando uma síntese analítica do trabalho desenvolvido, a avaliação crítica do cumprimento dos objetivos estipulados, o reconhecimento das limitações identificadas durante a pesquisa e a proposição de direcionamentos para trabalhos futuros.

\section{Síntese do trabalho, contribuições e objetivos atingidos}
\label{sec:concl_sintese}
\label{sec:concl_objetivos_atingidos}

O presente Trabalho de Conclusão de Curso concebeu, implementou e avaliou experimentalmente uma solução computacional de ponta a ponta para a extração, saneamento e classificação automática de transações financeiras em faturas de cartão de crédito no contexto brasileiro. Suas contribuições, confrontadas com os objetivos específicos estabelecidos na Introdução, são as seguintes:
\begin{enumerate}
	\item \textbf{Esteira de extração com preservação de leiaute} (objetivo 1): serviço de extração baseado em IBM Docling V2 que, em faturas nativas digitais, reduz a taxa de erro de caracteres de 52,82\% para 5,72\% frente à rasterização seguida de OCR (Tabela~\ref{tab:resultados_ocr_consolidados}). O componente de reconhecimento óptico local (Tesseract com filtro de confiança) foi implementado e exercitado na comparação, mas não acoplado ao serviço (Seção~\ref{sec:sol_extracao_hibrida}) -- objetivo atendido para documentos nativos digitais e parcialmente para documentos ópticos;
	\item \textbf{Saneamento léxico} (objetivo 2): rotina determinística de normalização de prefixos de intermediadores, ruídos alfanuméricos e aliases, cujo efeito mostrou-se \textbf{condicional à representação}: decisivo para a busca vetorial densa e neutro a levemente prejudicial para os classificadores esparsos (Seção~\ref{sec:val_ablacao_discussao}) -- objetivo atendido;
	\item \textbf{Linhas de base, representações densas e modelo no nível de byte} (objetivos 3 e 4): TF-IDF com Naive Bayes e Random Forest, busca vetorial densa sobre pgvector e ajuste fino do ByT5, todos calculados, comparados e documentados sob o mesmo teste cego -- objetivos atendidos;
	\item \textbf{Agente em grafo de estados} (objetivo 5): agente LangGraph que articula busca vetorial, raciocínio com LLM local, ferramenta de busca na web e interrupção para confirmação humana. O grafo foi construído e avaliado; dentro dele, o ramo vetorial foi o componente determinante e o ramo LLM, o gargalo de latência e de acurácia (Seção~\ref{sec:val_ablacao_discussao}). Três mecanismos previstos não chegaram a influenciar a decisão na versão avaliada -- a busca na web (exercitada em execução complementar comprometida por falhas de conexão), o histórico de \textit{marketplace} (recuperado, mas não injetado no \textit{prompt}) e a confirmação humana em tempo real (simulada, por limitação do \textit{checkpointer}) -- objetivo atendido em sua arquitetura e parcialmente em sua operação;
	\item \textbf{Infraestrutura em camadas e MLOps local} (objetivo 6): esteira conteinerizada de baixo custo apoiada em MinIO, Apache Kafka, PostgreSQL/pgvector, MLflow e Logto, com rastreamento de experimentos no nível de arquitetura e verificação empírica de necessidade de dados sob a LGPD. O registro por inferência, o controle de acesso por papéis nas APIs e a camada de consumo analítico permanecem como evolução (Seção~\ref{sec:concl_limitacoes}) -- objetivo atendido no essencial;
	\item \textbf{Validação sob protocolo rigoroso} (objetivo 7): cinco arquiteturas avaliadas sobre 1.660 transações reais em 11 categorias, com teste cego de 334 transações por estabelecimento disjunto, validação cruzada nas duas linhas de base esparsas, ablação, intervalos de confiança por \textit{bootstrap} e de Wilson, testes de McNemar e binomial, e rastreamento no MLflow -- objetivo atendido. A hipótese central foi \textbf{refutada} em suas duas subhipóteses para a esteira híbrida que a enuncia (acurácia de 61,4\% e F1-Macro de 60,0\%); entre as arquiteturas de referência, apenas o Naive Bayes supera o limiar de acurácia global, de forma marginal, e nenhuma supera o de F1-Macro (Seção~\ref{sec:val_respostas_questoes});
	\item \textbf{Achado metodológico transferível:} a quantificação do efeito de memorização de estabelecimentos sob validação cruzada convencional (Seção~\ref{sec:val_resultados_classificacao}), desenvolvida na conclusão 2 da Seção~\ref{sec:concl_principais_conclusoes};
	\item \textbf{Publicação e reprodutibilidade:} todo o código-fonte -- microsserviços, processos trabalhadores, agente e a esteira experimental numerada que reproduz os resultados do Capítulo~\ref{Cap:Ava_Exp} -- está publicado em repositório Git público (Apêndice~\ref{apendice:repositorio}), com o acervo de faturas reais excluído da distribuição em observância à LGPD. A revisão final do código frente ao texto, incorporada a esta versão, explicitou as divergências entre o que foi projetado e o que foi efetivamente exercitado, registradas nos capítulos correspondentes.
\end{enumerate}

\section{Principais conclusões da pesquisa}
\label{sec:concl_principais_conclusoes}

Da evidência apresentada no Capítulo~\ref{Cap:Ava_Exp}, destacam-se as seguintes conclusões, cada uma remetida à questão de pesquisa que a origina (Seção~\ref{sec:sol_questoes_pesquisa}):

\begin{enumerate}
	\item \textbf{Em faturas nativas digitais, preservar a topologia tabular na ingestão é condição necessária para a qualidade da classificação subsequente} (Q3, quanto ao erro de leitura óptica). Descartar a camada de texto do PDF e rasterizá-lo para reconhecimento óptico faz mais da metade dos caracteres da descrição chegar corrompida ao classificador (CER de 52,82\% contra 5,72\%; casamento exato de 3,3\% contra 86,2\%), e nenhum saneamento posterior recupera essa perda. A decisão de projeto mais consequente da esteira está na ingestão, não na modelagem;

	\item \textbf{O protocolo de avaliação é determinante para a leitura dos resultados -- a contribuição metodológica mais relevante desta pesquisa} (Q2). Para as duas arquiteturas em que a comparação foi mensurada, a diferença entre validação cruzada convencional e teste cego por estabelecimentos disjuntos foi de 23 a 25 pontos percentuais de F1-Macro, evidenciando que parcela substancial do desempenho usualmente reportado pode ser memorização de comércios já observados, e não generalização;

	\item \textbf{Em regime de baixo volume rotulado, uma representação esparsa por $n$-gramas de caractere com classificador clássico iguala o desempenho de \textit{embeddings} densos multilíngues a uma fração do custo, e supera com folga o ajuste fino de um modelo no nível de byte e um agente apoiado em LLM local -- a contribuição empírica mais relevante do trabalho} (Q1 e Q4). Comparadas em sua melhor configuração, as duas famílias de representação são estatisticamente indistinguíveis (69,3\% contra 66,8\% de F1-Macro; McNemar com $p = 0{,}86$), ao passo que o ByT5 (48,3\%) e o agente (60,0\%) ficam claramente abaixo. Sofisticação arquitetural não compensa escassez de dados rotulados; o que separa as duas primeiras abordagens neste regime é o custo, e o benefício do saneamento léxico é condicional à representação, não uma boa prática universal (Tabela~\ref{tab:ablacao});

	\item \textbf{No agente híbrido, o gargalo está no raciocínio, não na recuperação -- e a configuração avaliada do ramo de raciocínio não o testou em condições favoráveis} (Q3 e Q5). O ramo vetorial resolveu 98\% das transações que recebeu, em dezenas de milissegundos; o ramo LLM local acertou 21\% delas, a cerca de 14 segundos por transação, sob uma diretriz antialucinação que se converteu em abstenção sistemática para a classe residual. Quatro fatores confundem a atribuição dessa diferença ao modelo: o \textit{prompt} não calibrado (Apêndice~\ref{apendice:system_prompt}), a ausência dos vizinhos recuperados como exemplos no \textit{prompt} (o ramo operou em \textit{zero-shot}, e não como geração aumentada por recuperação), o metadado de origem fixo e a indisponibilidade efetiva da ferramenta de busca. Nenhum deles foi isolado experimentalmente; a leitura mais alinhada à literatura de agentes é que um modelo de 4 bilhões de parâmetros efetivos, sem evidência recuperada no \textit{prompt}, é insuficiente como cérebro do agente;

	\item \textbf{A proveniência do rótulo é um fator de desempenho e deve ser tratada como variável de estratificação obrigatória} (Q2). O F1-Macro do melhor modelo cai de 82,4\% para 39,4\% quando se passa do subconjunto rotulado por dicionário léxico para o rotulado pela categoria nativa do emissor, e a queda ocorre nas cinco arquiteturas (Tabela~\ref{tab:circularidade}). Reportar uma métrica agregada sobre acervos de rotulagem heterogênea oculta variação dessa magnitude;

	\item \textbf{Em taxonomias de despesa com cauda longa, acurácia global e F1-Macro medem propriedades distintas e podem apontar em direções opostas sobre o mesmo recorte} (Q2). No subconjunto do Nubank, quatro arquiteturas atingem 98,0\% de acurácia e 36,2\% de F1-Macro simultaneamente (Seção~\ref{sec:val_desempenho_emissor}). A métrica primária deve ser declarada e justificada antes da avaliação, como feito na Seção~\ref{sec:val_protocolo}; e a verificação de uma hipótese sobre acurácia deve vir acompanhada de intervalo de confiança, sem o qual o resultado marginal do Naive Bayes (88,3\%, IC 95\% [84,4\%; 91,3\%]) seria lido como sustentação inequívoca.
\end{enumerate}

\textbf{Recomendação prática.} Para quem precise categorizar faturas de cartão em um acervo pessoal de porte semelhante, a evidência deste trabalho recomenda: um classificador esparso (TF-IDF com $n$-gramas de caractere e Random Forest) como primeira linha, pelo desempenho e pelo custo desprezível; busca vetorial densa para estabelecimentos já observados, em que sua acurácia condicional é próxima de 100\%; modelo de linguagem restrito aos casos residuais, com evidência recuperada no \textit{prompt} e confirmação humana obrigatória; e ajuste fino de modelos no nível de byte apenas com acervos rotulados uma ordem de grandeza maiores.

Em síntese, esta monografia entregou uma esteira completa, reprodutível e executável em infraestrutura local, submeteu-a a um protocolo de avaliação mais severo que o usual na área e documentou, com a mesma precisão, o que foi demonstrado, o que foi apenas projetado e o que resta investigar.

\section{Limitações do trabalho}
\label{sec:concl_limitacoes}

Além da refutação da hipótese de desempenho preditivo, apontam-se as seguintes limitações, cada uma discutida na seção indicada:
\begin{enumerate}
	\item \textbf{Ferramenta de busca na web:} exercitada apenas em execução complementar em que os \textit{backends} de busca falharam por conexão em 14\% das transações; o resultado principal não a inclui (Seção~\ref{sec:val_ambiente});
	\item \textbf{Metadado de origem fixo no teste cego:} o LLM recebeu origem incorreta para todo o subconjunto do Nubank, e a componente ``origem do arquivo'' da Q4 não foi variada (Seções~\ref{sec:val_ambiente} e~\ref{sec:val_respostas_questoes});
	\item \textbf{Histórico de \textit{marketplace} e detecção de Mercado Livre:} o histórico recuperado não é injetado no \textit{prompt}, e um defeito de chave impede o roteamento do \textit{marketplace} mais frequente pelo ramo de histórico (Seção~\ref{sec:sol_agente_inteligente}; Apêndice~\ref{apendice:normalizacao});
	\item \textbf{HITL simulado nas duas modalidades:} a interrupção do grafo é retomada programaticamente tanto no teste cego quanto no modo de gravação; a espera por um revisor em tempo real exige um \textit{checkpointer} persistente e uma interface de revisão não implementadas (Seções~\ref{sec:sol_agente_inteligente} e~\ref{sec:val_mlops});
	\item \textbf{Ramo de raciocínio sem evidência recuperada:} o LLM decide em \textit{zero-shot}; a ausência dos vizinhos no \textit{prompt} é uma causa candidata, não isolada, da acurácia condicional de 21\% (Seção~\ref{sec:val_ablacao_discussao});
	\item \textbf{Ancoramento de rótulos:} o ancoramento interno ao nó de estruturação não é distinguível a posteriori do rótulo bruto do modelo, pois este não foi persistido; o ancoramento externo, verificável, não alterou nenhuma predição (Seção~\ref{sec:val_resultados_classificacao});
	\item \textbf{Extração óptica não acoplada nem validada:} o componente Tesseract não está integrado ao serviço de extração, e a comparação de OCR foi conduzida sobre PDFs nativos rasterizados, não sobre documentos genuinamente escaneados (Seções~\ref{sec:sol_extracao_hibrida} e~\ref{sec:val_resultados_ocr});
	\item \textbf{Porte, rotulagem e circularidade do acervo:} 1.660 transações de dois portadores, rotuladas semiautomaticamente, com 11 das 14 categorias povoadas e classes de cauda com 4 a 5 estabelecimentos; o F1-Macro cai sistematicamente no subconjunto rotulado pelo emissor, sem que se possa distinguir circularidade de ruído de rótulo (Seções~\ref{sec:val_dataset_classificacao} e~\ref{sec:val_circularidade});
	\item \textbf{Validação cruzada restrita e não agrupada:} a diferença entre validação cruzada e teste cego foi medida apenas para as duas linhas de base esparsas, e a validação cruzada empregada não agrupa por estabelecimento; uma validação cruzada agrupada (\texttt{StratifiedGroupKFold}) confirmaria o achado dentro do próprio treino, com cinco estimativas (Seção~\ref{sec:val_resultados_classificacao});
	\item \textbf{Desbalanceamento sem tratamento sistemático:} apenas o Random Forest recebeu ponderação de classes; a variante Complementar do Naive Bayes e técnicas de reamostragem não foram avaliadas (Seção~\ref{sec:fund_metricas_classificacao});
	\item \textbf{Hiperparâmetros e limiares não sintonizados:} o conjunto de validação foi incorporado ao treino; hiperparâmetros e os três limiares de roteamento foram fixados \textit{a priori} (Seções~\ref{sec:val_protocolo} e~\ref{sec:sol_agente_inteligente});
	\item \textbf{Reprodutibilidade do ramo LLM:} a configuração do modelo e do servidor foi documentada retroativamente a partir dos registros do LM Studio, e não pela esteira experimental no momento da execução; a reprodutibilidade bit a bit permanece sujeita a efeitos numéricos do motor de inferência (Seção~\ref{sec:val_ambiente});
	\item \textbf{Rastreabilidade, controle de acesso e testes:} o registro por inferência (\texttt{log\_classificacoes}), a validação de \textit{tokens} nas APIs, o isolamento multi-inquilino e a suíte de testes unitários não foram implementados; a aplicação \textit{web} é um protótipo sem painel orçamentário nem tela de confirmação (Seções~\ref{sec:sol_arquitetura}, \ref{sec:sol_lgpd} e~\ref{sec:sol_testes});
	\item \textbf{Modelo de \textit{embeddings} genérico:} o modelo multilíngue de paráfrase foi usado sem ajuste fino ao domínio de nomes de estabelecimento (Seção~\ref{sec:val_resultados_classificacao});
	\item \textbf{Levantamento bibliográfico:} cobertura mais rasa da string terciária, triada pelos 50 registros mais citados de cada base (Apêndice~\ref{apendice:protocolo_busca});
	\item \textbf{Desvio estatístico não mensurado:} \textit{Data Drift} e \textit{Concept Drift} motivam o ciclo de MLOps, mas não foram caracterizados, embora o acervo temporalmente ordenado o permitisse (Seção~\ref{sec:val_respostas_questoes}).
\end{enumerate}

\section{Trabalhos futuros}
\label{sec:concl_trabalhos_futuros}

Os desdobramentos organizam-se em dois blocos. O primeiro, de \textbf{consolidação dos resultados atuais}, reúne execuções de baixo custo sobre a esteira existente, cada uma respondendo a uma limitação da Seção~\ref{sec:concl_limitacoes}:
\begin{enumerate}
	\item \textbf{Validação com rótulos independentes} (limitação 8, de maior prioridade): constituir um subconjunto de teste validado manualmente, sem acesso ao dicionário léxico nem à categoria do emissor, para quantificar e descontar o efeito de circularidade; complementarmente, aplicar \textit{Confident Learning} para priorizar rótulos suspeitos à revisão;
	\item \textbf{Reexecução do ramo LLM em condições controladas} (limitações 1 a 5): sobre as mesmas transações roteadas ao LLM no teste cego, reexecutar o ramo com (a) os três vizinhos recuperados injetados no \textit{prompt} como exemplos rotulados -- convertendo o ramo em geração aumentada por recuperação --, (b) o metadado de origem correto, (c) o \textit{prompt} corrigido (regra 5 restaurada, campos revisados), (d) a ferramenta de busca com \textit{backend} estável e (e) um modelo local de maior porte (12 a 27 bilhões de parâmetros), isolando o efeito de cada fator na acurácia condicional de 21\%. Persistir o rótulo bruto do modelo para tornar o ancoramento auditável;
	\item \textbf{Validação cruzada agrupada, balanceamento e sintonia} (limitações 9 a 11): repetir a validação cruzada com \texttt{StratifiedGroupKFold} por estabelecimento; avaliar \textit{Complement Naive Bayes}, ponderação de classes e reamostragem; e empregar o conjunto de validação para sintonizar hiperparâmetros e os três limiares de roteamento;
	\item \textbf{Ajuste fino contrastivo do modelo de \textit{embeddings}} (limitação 14): treinar o codificador de sentenças com pares de descrições do mesmo estabelecimento ou categoria a partir do conjunto de treino;
	\item \textbf{Correções de integração} (limitações 3, 4, 7 e 13): corrigir a chave de detecção de \textit{marketplace} e injetar o histórico no \textit{prompt}; adotar um \textit{checkpointer} persistente e uma tela de confirmação na aplicação \textit{web}; acoplar o Tesseract ao serviço de extração como \textit{fallback}; alimentar \texttt{log\_classificacoes}; validar os \textit{tokens} do Logto nas APIs; e escrever a suíte de testes unitários das funções determinísticas;
	\item \textbf{Bateria de OCR sobre documentos genuinamente escaneados} (limitação 7), estendendo-a também às 49 competências pareadas do acervo.
\end{enumerate}

O segundo bloco, de \textbf{expansão}, reúne desenvolvimentos de maior alcance: integração com as APIs do Open Finance Brasil para ingestão contínua de extratos; arquitetura multiagente com especialistas em conciliação orçamentária, detecção de fraudes e consultoria tributária; modelos de visão e linguagem ultracompactos em dispositivo móvel para extração e categorização em borda; aprendizado federado para compartilhar, com preservação de privacidade, o conhecimento sobre novos estabelecimentos entre usuários; materialização das camadas Silver e Gold em formato de tabela aberto com transações ACID, conectadas a um motor SQL distribuído, a painéis gerenciais e a um orquestrador de fluxos, completando as propriedades de \textit{lakehouse}; e caracterização empírica de desvios de dados e de conceito sobre o acervo temporalmente ordenado.
